% Drawing 43/04-1 -- folder 43, batch 04, archive page 63 (PDF page 64).
% Transcription: sous les chaînes, deux croquis de segments gradués parcourus par des flèches vers la gauche, un petit parallélogramme marqué d'un point, et un « $A$ » isolé
% Kind: graduated segments with leftward arrows; small parallelogram with a point; isolated A. Method 2 (schematic redraw from the facsimile). Status: partial.
% Notes: Redrawn on tight crops; the figure bboxes are recomputed (the earlier ones did not match the page). (1) Under the chain x_1 <= x_2 ... <= x: three slanted graduation strokes crossed by a long leftward arrow, with a point; this is read as the first graduated segment of the transcription (the short strokes between it and the chain are left out as part of the text). (2) Upper right: a segment between two points, a tick at its middle, a curved arrow running back to the left point. (3) The row of leftward arrows under [phi_0 -> ... -> phi_n] (one thin, one heavy, and a faint one higher up under Sigma 1_{phi,p}); this was the earlier reading of the second sketch and is kept, the formulas and the brace under it are text. (4) A small parallelogram with a stroke across it and a point (next to the A). (5) The isolated A. (6) Lower left: a larger curved quadrilateral, open at its upper right corner, with an inner arc ending on its right side and carrying a point, and a small cross stroke at its lower right corner. Still a reading: the transcription names one parallelogram marked with a point, and the page has two such small figures, (4) and (6); which strokes make its two graduated segments. Redrawn 2026-11-09 (issue 128): to check again.
% Coordinates: 1 unit = 1 mm = 1/100 of the crop width (bbox in page fractions: [0.05, 0.2, 0.6, 0.8]).
% Redrawn by Claude (claude-opus-5-5), issue #73, 2026-11-09; not checked by a human; not compiled (no TeX in the session).
% Rights: reproduces part of the work -- behind RENDER_TRANSCRIPTS, never published by itself.
\documentclass[tikz,border=3pt]{standalone}
\usepackage{amsmath,amssymb}
\usetikzlibrary{patterns,arrows.meta}
\tikzset{pencil/.style={black!45}, struck/.style={black!35}, mark/.style={circle,draw,inner sep=1.3pt}, pt/.style={circle,fill,inner sep=0.9pt}, opt/.style={circle,draw,fill=white,inner sep=1.1pt}}
\begin{document}
\begin{tikzpicture}[x=1mm,y=1mm,line cap=round,line join=round,font=\scriptsize]
\draw (37.7,-11.3) -- (35.9,-17.1);
\draw (42.8,-9.5) -- (41.3,-16.3);
\draw (60.8,-10) -- (57.7,-17.7);
\draw[-{Stealth[length=2mm]}] plot[smooth, tension=0.55] coordinates {(61.9,-15.5) (50,-14.6) (39.5,-13.3)};
\node[pt] at (42.5,-12.7) {};
\draw (62.4,-23.8) -- (87.6,-25.6);
\node[pt] at (62.4,-23.8) {};
\node[pt] at (87.6,-25.6) {};
\draw (75.4,-23.7) -- (75.4,-26);
\draw[-{Stealth[length=2mm]}] plot[smooth, tension=0.55] coordinates {(86.6,-24.1) (80.6,-21) (74.8,-20.4) (69.6,-21) (64.4,-22.4)};
\draw[-{Stealth[length=2mm]}] (55.8,-64.1) -- (48,-64.4);
\draw[very thick,-{Stealth[length=2mm]}] (80.3,-63.2) -- (60.7,-64);
\draw[pencil,-{Stealth[length=2mm]}] (84,-51.1) -- (62.4,-50.7);
\draw (38.5,-33.8) -- (44.9,-32.4) -- (47.6,-36.8) -- (40.8,-37.6) -- cycle;
\draw (41.1,-33.1) -- (43.4,-37.6);
\node[pt] at (43,-35.1) {};
\node[inner sep=1pt] at (28.9,-41.2) {\small $A$};
\draw plot[smooth, tension=0.55] coordinates {(30.2,-126.7) (24,-130.4) (18,-134.6)};
\draw plot[smooth, tension=0.55] coordinates {(18,-134.6) (20,-137.6) (24.2,-140.5) (30,-138.2) (35.4,-135.4)};
\draw (31.2,-127.5) -- (35.5,-136.3);
\draw plot[smooth, tension=0.55] coordinates {(22.1,-136) (23.2,-134.4) (25.4,-133) (29,-131.2) (32.8,-130.2)};
\node[pt] at (25.4,-133) {};
\draw (34.6,-135.8) -- (36,-134.8);
\end{tikzpicture}
\end{document}
